\documentclass[../cours_spe_trigo.tex]{subfiles}
\begin{document}
\setcounter{chapter}{17}
\chapter{Méthodes, pièges et carte des démonstrations}
\section{Reconnaître ce qui est une définition et ce qui est démontré}
\begin{center}
\begin{tabularx}{\linewidth}{@{}p{.28\linewidth} X@{}}
\toprule
\textbf{Statut} & \textbf{Exemples et justification}\\
\midrule
Définition & $\pi=C/D$ après preuve que ce quotient est constant ; mesure en radians $\ell/R$ ; sinus et cosinus comme coordonnées du cercle.\\
Théorème de géométrie & Pythagore ; loi des cosinus ; loi des sinus ; produit scalaire et angle.\\
Conséquence géométrique & $C=2\pi R$ ; $180^\circ=\pi$ rad ; $\sin^2x+\cos^2x=1$ ; symétries et périodicité.\\
Conséquence algébrique & Formules d'addition, angles doubles, produit--somme, somme--produit, équations.\\
Résultat d'analyse & Limite $\sin x/x\to1$, dérivées, primitives, approximations de Taylor.\\
\bottomrule
\end{tabularx}
\end{center}
\section{Les erreurs fréquentes et leurs remèdes}
\begin{enumerate}[label=\arabic*.]
\item \textbf{Confondre $2\pi$ et $2\pi R$.} $2\pi$ est un angle (en radians), mais $2\pi R$ est une longueur.
\item \textbf{Confondre longueur et aire.} $C=2\pi R$ et $\mathcal A=\pi R^2$, avec unités différentes.
\item \textbf{Utiliser les degrés dans une formule d'analyse.} Pour $\sin x/x\to1$, les radians sont indispensables.
\item \textbf{Dire $\sqrt{\sin^2x}=\sin x$.} Le résultat général est $\sqrt{\sin^2x}=|\sin x|$.
\item \textbf{Écrire $\sin(a+b)=\sin a+\sin b$.} C'est faux en général : utiliser la formule d'addition.
\item \textbf{Diviser par une expression qui peut être nulle.} Dans une équation, cela peut supprimer des solutions.
\item \textbf{Oublier $k\in\mathbb Z$.} Les solutions générales se répètent à chaque période.
\item \textbf{Confondre $\arcsin(\sin x)$ et $x$.} L'égalité n'est vraie que dans l'intervalle principal choisi.
\item \textbf{Ignorer le domaine de la tangente.} Elle n'existe pas lorsque $\cos x=0$.
\end{enumerate}
\section{Chaîne logique du cours}
\begin{center}
\begin{tikzpicture}[node distance=6mm and 8mm,
nodeStep/.style={draw=bleu,fill=bleuPale,rounded corners=2mm,text width=3.7cm,align=center,inner sep=6pt,font=\small},
arr/.style={-Latex,thick,bleu}]
\node[nodeStep] (a) {Similitude, Pythagore, cercle et disque};
\node[nodeStep,below=of a] (b) {$\pi=C/D$ et $C=2\pi R$};
\node[nodeStep,below=of b] (c) {Radians, arcs et secteurs};
\node[nodeStep,below=of c] (d) {Cercle unité et coordonnées};
\node[nodeStep,below left=of d] (e) {Symétries, valeurs, périodicité};
\node[nodeStep,below right=of d] (f) {Rotations et formules d'addition};
\node[nodeStep,below=of $(e)!0.5!(f)$] (g) {Équations, fonctions, dérivées, applications};
\draw[arr] (a)--(b);\draw[arr] (b)--(c);\draw[arr] (c)--(d);
\draw[arr] (d)--(e);\draw[arr] (d)--(f);
\draw[arr] (e)--(g);\draw[arr] (f)--(g);
\end{tikzpicture}
\end{center}
Cette carte donne une stratégie de révision : si une identité paraît mystérieuse, \emph{remonter} d'un niveau, vers la rotation, la symétrie ou le triangle rectangle.
\end{document}
